\documentclass[12pt,a4paper]{article}

\usepackage{amsmath,amssymb,bm}
\usepackage[utf8]{inputenc}
\usepackage[T1]{fontenc}
\usepackage{geometry}
\usepackage{booktabs}
\usepackage{enumitem}
\usepackage{hyperref}
\usepackage{xcolor}
\usepackage{tabularx}
\usepackage{array}
\usepackage{parskip}
\usepackage[most]{tcolorbox}
\usepackage{setspace}

% ---- pseudocode box identical in spirit to the FOIL-algorithm image ----
\newtcolorbox{algobox}[1][]{%
  colback=white, colframe=black,
  boxrule=0.6pt, arc=0pt,
  left=6pt, right=6pt, top=4pt, bottom=4pt,
  before upper={\setstretch{1.25}\small},
  #1}

% Convenience macros for pseudocode typography
\newcommand{\kw}[1]{\textbf{#1}}          % keyword
\newcommand{\fn}[1]{\textsc{#1}}          % function name (small caps)
\newcommand{\var}[1]{\textit{#1}}         % variable name
\newcommand{\cm}[1]{\textcolor{gray}{#1}} % comment
\newcommand{\ind}{\hspace*{1.4em}}        % one indent level
\newcommand{\indd}{\hspace*{2.8em}}       % two indent levels
\newcommand{\inddd}{\hspace*{4.2em}}      % three indent levels
\newcommand{\assign}{\ensuremath{\leftarrow}}%

\geometry{margin=2.5cm}

\hypersetup{
    colorlinks=true,
    linkcolor=blue,
    citecolor=blue,
    urlcolor=blue
}

\definecolor{mplanned}{rgb}{0.0,0.45,0.70}
\definecolor{murgent}{rgb}{0.80,0.10,0.10}
\definecolor{mminor}{rgb}{0.10,0.50,0.10}

\newcommand{\planned}[1]{{\color{mplanned}#1}}
\newcommand{\urgent}[1]{{\color{murgent}#1}}
\newcommand{\minor}[1]{{\color{mminor}#1}}

\title{\textbf{Revision Plan}\\[4pt]
\large Manuscript SAMO-D-26-00346\\
\textit{``Frequency maximization of planar structures using\\quasi-static approximation in topology optimization''}}

\author{Piotr Tauzowski \and Andrzej Szczepańczyk \and Bartłomiej Błachowski}

\date{May 2026}

\begin{document}

\maketitle

% ============================================================
\section{Algorithm Pseudocodes and Comparison-Evidence Policy}
\label{sec:algorithms}
% ============================================================

This section presents the three algorithmic ideas in a uniform notation.
Algorithm~1 is a literature-level synopsis of the published Du--Olhoff
bound formulation; it is not a claim that any repository implementation is
paper-faithful or reproduces the published benchmark.  The only Olhoff-based
implementation admitted to production comparisons is
\texttt{analysis/OlhoffApproach}, which must be described as the
\emph{local Olhoff-inspired comparison implementation}.  The unsuccessful
\texttt{OlhoffApproachExact} reconstruction campaign is archived diagnostic
work and is not part of the reviewer-evidence chain.
Notation throughout:
$N_e$ = number of finite elements;
$n$ = index of the lowest targeted eigenfrequency (the one being maximized);
$J$ = number of candidate eigenpairs tracked ($J \geq n$);
$p,d$ = SIMP penalization exponents for stiffness and mass, respectively;
$x_e \in [x_{\min},1]$ = element design (density) variable;
$\rho_e$ = physical density after filtering/projection
  (in the sensitivity expressions $(\cdot)'_{\rho_e}$ denotes
  $\partial(\cdot)/\partial\rho_e$; for unfiltered formulations $\rho_e = x_e$);
$E_0, E_{\min}$ = Young's modulus of solid and void phases in the SIMP model
  ($E(\rho_e)=E_{\min}+\rho_e^p(E_0-E_{\min})$);
$\mathbf{K}_e^0$ = element stiffness matrix evaluated at unit density ($\rho_e=1$);
$c = \mathbf{f}^\top\mathbf{u}$ = structural compliance (objective in static sub-problems);
$\delta_{jk}$ = Kronecker delta (used in mass-orthonormality $\boldsymbol{\varphi}_j^\top\mathbf{M}\boldsymbol{\varphi}_k=\delta_{jk}$);
$V_0$ = total domain volume;
$V_e$ = volume of element $e$;
$V^* = V_f\cdot V_0$ = prescribed material volume;
$r_{\min}$ = sensitivity filter radius;
$\varepsilon$ = generic outer convergence tolerance;
$\varepsilon_1,\varepsilon_2$ = Stage~1 and Stage~2 convergence tolerances (Yuksel--Yilmaz);
$\varepsilon_{\mathrm{inner}}$ = inner MMA loop tolerance (Olhoff--Du);
$\varepsilon_{\mathrm{tol}}$ = main loop tolerance (Proposed).

% -----------------------------------------------------------
\subsection{Algorithm 1: Published Du--Olhoff Bound Formulation}
\label{subsec:alg-olhoff}
% -----------------------------------------------------------

At the literature level, the method maximizes a scalar bound variable
$\beta \leq \omega_j^2$, $j = n,\ldots,J$, through a nested double-loop
procedure.  The outer loop solves the eigenvalue problem and assembles
sensitivities; the inner loop runs MMA iteratively until the design
increments converge \cite{Olhoff2014}.  The pseudocode below records that
published conceptual structure only.  It must not be read as a validated
mapping to \texttt{OlhoffApproach}, and it supplies no empirical benchmark
for the revision.

\medskip
\begin{algobox}
\kw{function} \fn{Olhoff-Du}(\var{mesh}, \var{J}, \var{$V_f$})
\kw{returns} optimal density field $x^*$\\[2pt]
\ind\kw{inputs}: \var{mesh}, finite element mesh and boundary conditions\\
\ind\hspace*{4.2em}\var{J}, number of candidate eigenfrequencies ($J \geq n$)\\
\ind\hspace*{4.2em}$V_f$, prescribed volume fraction\\[2pt]
\ind\kw{local variables}:
  $x = \{x_e\}$, design variables, initially $x_e = V_f$\\
\ind\hspace*{9.5em}
  $\beta$, scalar bound variable\\
\ind\hspace*{9.5em}
  $\lambda_j, \omega_j, \boldsymbol{\varphi}_j$,
  eigenvalues, frequencies, eigenvectors, $j=1,\ldots,J$\\
\ind\hspace*{9.5em}
  $N$, detected multiplicity of $\omega_n$\\
\ind\hspace*{9.5em}
  $\tilde{\lambda}$, representative eigenvalue of the degenerate cluster
    (e.g.\ mean of $\lambda_n,\ldots,\lambda_{n+N-1}$), used in the
    generalized gradient formula\\
\ind\hspace*{9.5em}
  $\mathbf{f}_{sk}$, generalized gradient vectors (or
  $\nabla\!\lambda_j$ if $N=1$)\\
\ind\hspace*{9.5em}
  $\Delta\var{x}$, vector of design variable increments\\[6pt]
%
$x_e \assign V_f$ \quad for all $e = 1,\ldots,N_e$;
\quad $\beta \assign 0$\\[4pt]
%
\kw{repeat} \cm{// outer (main) loop}\\
\ind\cm{// Step 1 — Eigensolve}\\
\ind solve $\mathbf{K}(\var{x})\,\boldsymbol{\varphi}_j
           = \omega_j^2\,\mathbf{M}(\var{x})\,\boldsymbol{\varphi}_j$,
      $\boldsymbol{\varphi}_j^{\top}\mathbf{M}\boldsymbol{\varphi}_k=\delta_{jk}$,
      $j=1,\ldots,J$\\
\ind detect multiplicity $N$ of $\omega_n$\\[3pt]
%
\ind\cm{// Step 2 — Sensitivity analysis}\\
\ind\kw{if} $N > 1$ \kw{then}\\
\indd compute generalized gradients
      $\mathbf{f}_{sk}=\{\boldsymbol{\varphi}_s^{\top}
      (\mathbf{K}'_{\rho_e}-\tilde{\lambda}\mathbf{M}'_{\rho_e})
      \boldsymbol{\varphi}_k\}^{\top}$,\;
      $s,k=n,\ldots,n{+}N{-}1$\\
\ind\kw{else}\\
\indd compute standard gradients
      $(\lambda_j)'_{\rho_e}=
      \boldsymbol{\varphi}_j^{\top}
      (\mathbf{K}'_{\rho_e}-\lambda_j\mathbf{M}'_{\rho_e})
      \boldsymbol{\varphi}_j$,\;
      $j=n,\ldots,J$\\[3pt]
%
\ind\cm{// Step 3 — Inner loop: MMA sub-problem for design increments}\\
\ind\kw{repeat} \cm{// inner loop}\\
\indd solve MMA sub-problem (bound form, Eq.~19 of \cite{Olhoff2014})
      for increments $\Delta x_e$,\; $e=1,\ldots,N_e$\\
\ind\kw{until} $\|\Delta\var{x}\| < \varepsilon_{\mathrm{inner}}$\\[3pt]
%
\ind\cm{// Step 4 — Update design variables}\\
\ind $x_e \assign x_e + \Delta x_e$ \quad for all $e$\\
\ind $\beta \assign \min_{j=n,\ldots,J}\omega_j^2(\var{x})$\\[4pt]
%
\kw{until} $\|\Delta\var{x}\| < \varepsilon$\\[4pt]
\kw{return} $x$
\end{algobox}

% -----------------------------------------------------------
\subsection{Algorithm 2: Yuksel--Yilmaz Two-Stage Static Method}
\label{subsec:alg-yuksel}
% -----------------------------------------------------------

Two sequential static compliance minimizations replace the per-iteration
eigensolve.
Stage~1 produces a static displacement field as the initial eigenvector
estimate; Stage~2 iteratively refines this estimate through a
design-dependent inertial load \cite{Yuksel2025}.

\medskip
\begin{algobox}
\kw{function} \fn{Yuksel-Yilmaz}(\var{mesh}, $V_f$,
$\mathbf{f}_s$) \kw{returns} optimal density field $x^*$\\[2pt]
\ind\kw{inputs}: \var{mesh}, finite element mesh and boundary conditions\\
\ind\hspace*{4.2em}$V_f$, prescribed volume fraction\\
\ind\hspace*{4.2em}$\mathbf{f}_s$, static point load
  (midpoint for symmetric BCs; from solid eigensolve for asymmetric BCs)\\[2pt]
\ind\kw{local variables}:
  $x=\{x_e\}$, design variables, initially $x_e = V_f$\\
\ind\hspace*{9.5em}
  $\mathbf{u}_s$, static displacement from Stage~1\\
\ind\hspace*{9.5em}
  $\mathbf{u}_d$, current mode-shape estimate;
  $\hat{\mathbf{u}}_d$, its unit vector\\
\ind\hspace*{9.5em}
  $\mathbf{f}_d$, design-dependent inertial load vector\\[6pt]
%
$x_e \assign V_f$ \quad for all $e$\\[4pt]
%
\cm{// Stage 1 — mode-shape estimation via static compliance}\\
\kw{repeat}\\
\ind solve $\mathbf{K}(\var{x})\,\mathbf{u}_s = \mathbf{f}_s$\\
\ind compute
    $\frac{\partial c}{\partial x_e}
    = p x_e^{p-1}(E_0-E_{\min})\,\mathbf{u}_{s,e}^{\top}\mathbf{K}_e^0\,\mathbf{u}_{s,e}$,
    \quad $e=1,\ldots,N_e$\\
\ind apply sensitivity filter; update $x$ by OC bisection on
    $V(\var{x})=V_f\cdot V_0$\\
\kw{until} $\max_e|\Delta x_e| < \varepsilon_1$\\[4pt]
%
$\mathbf{u}_d \assign \mathbf{u}_s$
\quad\cm{// initialise eigenvector estimate from Stage~1 result}\\[4pt]
%
\cm{// Stage 2 — self-consistent iterative compliance with inertial load}\\
\kw{repeat}\\
\ind $\hat{\mathbf{u}}_d \assign \mathbf{u}_d\,/\,\|\mathbf{u}_d\|$
  \quad\cm{// normalise current mode estimate}\\
\ind $\mathbf{f}_d \assign \mathbf{M}(\var{x})\,\hat{\mathbf{u}}_d$
  \quad\cm{// assemble design-dependent inertial load}\\
\ind solve $\mathbf{K}(\var{x})\,\mathbf{u}_d = \mathbf{f}_d$\\
\ind compute $\frac{\partial c}{\partial x_e}
    = p x_e^{p-1}(E_0-E_{\min})\,\mathbf{u}_{d,e}^{\top}\mathbf{K}_e^0\,\mathbf{u}_{d,e}$\\
\ind apply sensitivity filter; update $x$ by OC bisection on
    $V(\var{x})=V_f\cdot V_0$\\
\kw{until} $\max_e|\Delta x_e| < \varepsilon_2$\\[4pt]
%
apply Heaviside post-processing to obtain solid--void design\\
solve $\mathbf{K}(\var{x}^*)\boldsymbol{\varphi}
      = \omega_1^2\,\mathbf{M}(\var{x}^*)\boldsymbol{\varphi}$
\quad\cm{// final eigensolve to report $\omega_1$}\\[4pt]
\kw{return} $x^*$
\end{algobox}

% -----------------------------------------------------------
\subsection{Algorithm 3: Proposed Quasi-Static Approximation}
\label{subsec:alg-proposed}
% -----------------------------------------------------------

A single eigensolve on the declared reference domain freezes the eigenpair
permanently.  The pseudocode below assumes the physical
$(\omega_j^{(0)})^2\mathbf{M}\boldsymbol{\Phi}_j$ inertial-load definition;
if A0 is resolved by retaining the current \texttt{semi\_harmonic}
implementation, replace this formula consistently with the implemented
reference-load definition.
Every subsequent iteration is a standard compliance minimization step;
the multi-mode load-weighting parameter~$\alpha$ provides simultaneous
two-mode control without any structural change to the algorithm.

\medskip
\begin{algobox}
\kw{function} \fn{Quasi-Static-Freq}(\var{mesh}, $V_f$,
$n_{\mathrm{modes}}$, $\alpha$) \kw{returns} optimal density field $x^*$\\[2pt]
\ind\kw{inputs}: \var{mesh}, finite element mesh and boundary conditions\\
\ind\hspace*{4.2em}$V_f$, prescribed volume fraction\\
\ind\hspace*{4.2em}$n_{\mathrm{modes}}$, number of modes to target
  (1 or 2)\\
\ind\hspace*{4.2em}$\alpha \in [0,1]$, load-weighting parameter\\[2pt]
\ind\kw{local variables}:
  $x=\{x_e\}$, design variables, initially $x_e=V_f$\\
\ind\hspace*{9.5em}
  $(\omega_j^{(0)},\boldsymbol{\Phi}_j)$,
  reference eigenpairs $j=1,\ldots,n_{\mathrm{modes}}$,
  \emph{fixed throughout}\\
\ind\hspace*{9.5em}
  $\mathbf{f}_\alpha$, weighted inertial load vector\\
\ind\hspace*{9.5em}
  $\mathbf{u}$, static displacement field\\[6pt]
%
\cm{// Initialisation — single eigensolve on declared reference domain}\\
$x_e \assign V_f$ \quad for all $e = 1,\ldots,N_e$\\
assemble $\mathbf{K}_0$ and $\mathbf{M}_0$ at the chosen reference design\\
solve $\mathbf{K}_0\boldsymbol{\Phi}_j
      = (\omega_j^{(0)})^2\mathbf{M}_0\boldsymbol{\Phi}_j$
  \quad for $j=1,\ldots,n_{\mathrm{modes}}$
  \quad\cm{// performed \emph{once}, then $(\omega_j^{(0)},\boldsymbol{\Phi}_j)$ frozen}\\[4pt]
%
\cm{// Main loop — pure static compliance minimization, no eigensolve}\\
\kw{repeat}\\
\ind assemble $\mathbf{K}(\var{x})$ and $\mathbf{M}(\var{x})$\\
\ind $\mathbf{f}_\alpha \assign
  \alpha\,(\omega_1^{(0)})^2\mathbf{M}(\var{x})\boldsymbol{\Phi}_1
  + (1{-}\alpha)\,(\omega_2^{(0)})^2\mathbf{M}(\var{x})\boldsymbol{\Phi}_2$
  \quad\cm{// frozen $\boldsymbol{\Phi}_j$}\\
\ind solve $\mathbf{K}(\var{x})\,\mathbf{u} = \mathbf{f}_\alpha(\var{x})$\\
\ind compute
  $\dfrac{\partial c}{\partial x_e}
  = -p\,x_e^{p-1}(E_0-E_{\min})\,\mathbf{u}_e^{\top}\mathbf{K}_e^0\,\mathbf{u}_e$
  \quad\cm{// load sensitivity $\partial\mathbf{f}/\partial x_e$ set to zero}\\
\ind apply density-weighted sensitivity filter of radius $r_{\min}$\\
\ind update $x$ by OC bisection on volume constraint
  $\sum_e x_e V_e = V^*$\\
\kw{until} $\max_e|\Delta x_e| < \varepsilon_{\mathrm{tol}}$\\[4pt]
%
\cm{// Post-analysis}\\
solve $\mathbf{K}(\var{x}^*)\boldsymbol{\varphi}_j
      = \omega_j^2\,\mathbf{M}(\var{x}^*)\boldsymbol{\varphi}_j$,
  \;$j=1,\ldots,3$
  \quad\cm{// one final eigensolve to verify $\omega_1$}\\[4pt]
\kw{return} $x^*$
\end{algobox}

% -----------------------------------------------------------
\subsection{Performance-Comparison Scope}
\label{subsec:complexity}
% -----------------------------------------------------------

\subsubsection*{Evidence boundary}

The published Algorithm~1 and the local \texttt{OlhoffApproach} code play
different roles.  Algorithm~1 provides literature context for the bound
formulation.  The local code is an Olhoff-inspired implementation that may
be used only as a local empirical comparator.  Agreement between the two
has not been established, so observations from the local code cannot be
reported as properties, costs, convergence behaviour, or achieved optima of
the published Du--Olhoff method.

The attempted \texttt{OlhoffApproachExact} reconstruction did not establish
a paper-faithful implementation.  Its scripts, outputs, timings, iteration
histories, and terminal frequencies are archived diagnostics.  They must not
be used to infer a frequency gap to a Du--Olhoff optimum, a speedup relative
to a canonical implementation, a convergence advantage, or a mesh-scaling
law.

\subsubsection*{Permitted local comparison}

If an empirical performance comparison is retained in the revision, its
Olhoff-based column must be generated by and explicitly labelled
\emph{local \texttt{OlhoffApproach}}.  Such a comparison may describe only
the observed behaviour of the repository implementations under the stated
settings.  It must not use ``canonical'', ``exact'', ``reference'',
``paper-faithful'', or ``benchmark implementation'' as qualifiers for the
local comparator, and it must not reinterpret local endpoint differences as
gaps to the published optimum.

No quantitative timing, frequency-gap, convergence, memory, or scaling
table is carried forward in this planning document, and none is reinstated:
EXP1 and EXP5 are retired from the reviewer evidence chain
(\texttt{SCIENTIFIC\_DECISION\_EXP1\_EXP5.md}), so \textbf{no quantitative
cross-code performance benchmark remains anywhere in the revision}.  Any
separate reviewer-facing table or figure must be non-comparative --- it may
report the proposed method's own results, never a runtime, memory, speedup, or
scaling comparison against another implementation.  This policy removes the
unsupported empirical chain without requesting a replacement experiment.

\subsubsection*{Qualitative algorithmic distinction}

At the conceptual level, the published bound formulation repeatedly solves
an eigenvalue problem and uses eigenvalue sensitivities within its
optimization structure, whereas the proposed quasi-static approximation
freezes its reference eigenpair and performs static solves in its main loop.
This structural distinction motivates discussing computational work, but by
itself it does not establish a numerical speedup, convergence ranking, or
scaling advantage.  Those are empirical properties of specific
implementations and run conditions; for the revision, the only eligible
Olhoff-based implementation is the local \texttt{OlhoffApproach} comparator
defined above.
\bibliographystyle{plain}
\bibliography{../literature}

\end{document}
